\section{Related Work}
\label{sec:related}

\emph{Assurance cases and change impact.}
System change does not invalidate every assurance artifact.  It requires
change-impact analysis to identify which claims depend on premises that no
longer hold.  Kelly and McDermid give an impact analysis for deciding which
parts of an argument a change invalidates~\cite{kelly_safety_case_maintenance_2001}, and
dynamic safety cases carry it into operation so assurance is re-argued through
life~\cite{denney_dynamic_safety_cases_2015}.  That literature asks whether
engineers correctly re-argue safety after a known change.  We instantiate that
maintenance problem at two controller handoffs: extrapolating a bounded
raw-policy release test and transferring control to a fallback are separate
change-impact events with distinct premises.

\emph{Runtime assurance and safe controller updates.}
Safety filters and shields form a design space, and we position against it as
a whole~\cite{brunke_safe_learning_survey_2022}.
Simplex architectures retain a decision module and a baseline trusted under
the system fault model during online upgrades~\cite{seto_simplex_1998}; neural
and barrier-based Simplex add online retraining, reverse switching, and
switching conditions under bounded model error
\cite{phan_neural_simplex_2020,damare_bb_simplex_2025}; and MPSC and runtime
shields preserve recoverability or prevent unsafe actions
\cite{wabersich_linear_mpsc_2018,didier_adaptive_mpsc_2021,
alshiekh_shielding_2018}.  These lines already establish that safe switching
requires a recoverability or invariant-set argument.  Our controlled audit
deliberately evaluates a weaker mechanism---repeated lookahead for the outgoing
raw policy without a recovery certificate for the receiver---to measure the
consequence of treating those claims as interchangeable.  Policy repair, safe projection, joint
controller/certificate repair, robust conformal transfer, and shielded
sim-to-real transfer instead establish evidence for the deployment object
before use \cite{zhou_runtime_policy_repair_2020,lu_repair_preservation_2024,
chow_safe_policy_2021,yu_certificate_repair_2025,
mirzaeedodangeh_safe_updates_2026,hsu_sim_lab_real_2023}.

\emph{Shield retirement and attacked states.}
Könighofer et al.~\cite{konighofer_online_shielding_2022} establish the
non-adversarial half of this observation: policies trained under an online
shield inflict more unshielded violations than conventionally trained ones, so
shielding must be applied during execution.  They further note that a
finite-horizon shield needs sufficient lookahead, or the agent reaches states
in which every action is unsafe; the 41 empty kernels under our one-step action
filter instantiate that situation.  Our horizon sweep quantifies the tradeoff
between warning time and admission.  The duration study then shows why
raw-policy warning cannot replace the recoverability condition already required
by runtime-assurance designs: 19 paired failures occurred only under raw
execution, whereas seven occurred only after the lookahead switcher transferred
control to the LQR fallback.  Because the experiment evaluates no recovery
certificate, these outcomes do not locate the handoff states relative to a
recoverable set.  Carr et al.~\cite{carr_partial_shielding_2023} study a partially observable
bootstrap-to-retirement lifecycle and compare sudden with smooth shield
removal under trusted learning signals.  We reuse that published lifecycle as
a third-party case study and add bounded reward-record poisoning during the
protected bootstrap.  Carr et al. do not claim that a five-step test certifies
permanent execution; their workflow establishes the retirement setting, while
our paired experiment evaluates whether bootstrap-only data corruption can
alter the policy exposed at retirement.

Secure estimation characterizes when sensor attacks can be detected or
corrected~\cite{pasqualetti_attack_2013,fawzi_secure_2014}, and severe-attack
CBF filtering constructs plausible-state sets and safe-action kernels when the
state is not uniquely reconstructable~\cite{tan_safety_2024}; our model
consumes such a set and asks whether an adversarially updated action map stays
inside its certified kernel.  Related lines learn barriers or solve
adversarial reach--avoid problems under trusted learning signals
\cite{brunke_barrier_learning_2022,
capone_bandit_safe_control_2025,wu_certifying_2022,pham_certified_2024,
hsu_isaacs_2023,nguyen_gameplay_2024}.  Table~\ref{tab:delta} in
Appendix~\ref{sec:appendix-delta} places the closest work on the three axes
this paper combines.

\emph{Policy and reward poisoning.}
Policy poisoning has been studied in batch reinforcement learning and
control~\cite{ma_policy_poisoning_2019}; online attacks perturb rewards under
an $\ell_\infty$ budget and adapt to the
learner~\cite{zhang_reward_poisoning_2020}; extensions cover policy-gradient
and black-box deep learners, and targeted poisoning can teach a chosen policy
\cite{sun_online_rl_poisoning_2021,xu_deep_reward_poisoning_2025,
rakhsha_policy_teaching_2020,everitt_corrupted_reward_2017}.  None of this
line places the learner behind a runtime shield.

\emph{Backdoors on shielded and classical controllers.}
Jiang et al.\ implant backdoors in formal-language-guided \emph{safe} RL
policies, so that a trigger state reached at runtime produces a safety
violation~\cite{jiang_backdoor_safe_rl_cps_2024}.  This is the closest prior
composition, and the difference matters: their trigger must survive a shield
that is never retired, whereas our attack plants no input-triggered behavior.
Its effect becomes observable when protection is absent---after a bounded
release test in the controlled audit, or after shield retirement in the
third-party workflow.  The manipulated object is the learning record rather
than a deployment input.
BADControl poisons operational data in proportional-integral-derivative (PID)
and LQR loops to implant a vulnerability activated by a physical
trigger~\cite{burbano_badcontrol_2026}, which shares the trigger structure
without the learner or the shield.  Where online
reward-poisoning defenses model adaptive attacker--defender interaction
\cite{nika_reward_defense_2023}, we separate provenance, trusted
recomputation, detection, commit admission, and runtime containment.

Appendix~\ref{sec:appendix-support} applies classical homogenization and
second-order-cone feasibility \cite{schaible_fractional_duality_1976,
agrawal_dqcp_2020,bauschke_homogenization_2023,belloni_soc_feasibility_2008}
to the normalized image of a centered reward-to-go zonotope.

\section{Discussion and Limitations}

The clean cohort isolates horizon extrapolation: its five-step test directly
evaluates the final raw-policy snapshot without shield intervention, but does
not support a 120-step claim.  The duration study isolates a different premise:
warning about the outgoing policy does not establish LQR recovery from the
transfer state.  Reward-record poisoning stresses both contracts but creates
neither obligation.

In the Carr corpus, positive outcomes occurred only at low retirement-state
disagreement, but that region also contained non-positive outcomes.  The
53-run accounting (12 positive, 20 protective, and 21 unchanged) therefore
supports a conditional, not uniform, effect.  Trusted recomputation, commit
admission, and recovery-aware switching address different premises and incur
different coverage, latency, and intervention costs.

Adaptation should justify those costs by paired task benefit over freeze.  The
locked benign audit stopped after 0/12 improvements, so the evaluated settings
provide no such utility basis; this result does not extend to unevaluated
deployments.

\subsection{Limitations}
\label{sec:limits}

\emph{The marker is descriptive, not an attacker trigger.}  Retirement-state
disagreement is observed only after bootstrap.  Two historical REINFORCE
membership decisions also lack an executable deduplication rule, so we omit
the affected corpus statistic and counts.  The reproducible high-band result
is an obstacle-scoped exclusion condition, not a validated predictor; holding
out SAC identifies only 3 of 10 outcomes.

\emph{The formal results have bounded verification scope.}
The evidence record in Section~\ref{sec:threat} is dependency bookkeeping, not
a proof calculus, and yields no empirical result.  The linear parameter gate is
exact only on its finite abstraction; its continuous lift requires a justified
reachable-set cover, regularity and model-error bounds, and a robust margin.
The coverage audit does not establish continuous-state invariance.  Exact
reward influence covers one fixed standardized-advantage batch with global
gradient clipping, not coordinatewise clipping or multi-batch construction.

\emph{Empirical generality is bounded to Monte Carlo policy gradients.}  The
learned-baseline variant shares their return and attack channel, so it is not
an independent family.  The official Safe-Control-Gym PPO audit showed
checkpoint-level redirection without stable physical failure, and on the
third-party lifecycle PPO and SAC contained individual positive realizations
without reliable success across paired runs.

\emph{The horizon and workload claims remain experimental.}  The five-step
horizon was development-fixed, and its sweep uses the poisoned V3 cohort, not
the 20-run clean cohort.  We therefore do not estimate clean run-level failure
prevalence by horizon.  Both platforms expose recomputable deterministic
rewards; delayed labels and external KPIs motivate, but do not empirically
validate, the non-reconstructable-reward setting.

\emph{We assume access, and do not estimate its prevalence.}  The attack reads
learner state in addition to writing one scalar field
(Table~\ref{tab:trust-boundary}).  Predicate-aware edits evade the task-aware
checks; trusted recomputation closes the channel but falls outside the stated
applicability condition.  At 48 batches, switcher failures occur in two of five
controllers.  We neither certify an LQR recovery set nor identify safe handoff
states, limiting claims about prevalence and alternate switching times.
